\subsection{Influence AA from completion and elementary certificate bounds}
\label{ian:influence-form-uniform-aa}

Status: proved, for the uniform law and every finite product law, uniformly under positive-probability pinning, with exponents \((\alpha,\beta)=(12,12)\) and explicit constant \(c_\star=9/(2^{61}\pi^{12})\) (Theorem~\ref{ian:uniform-influence-aa-proved}, Corollary~\ref{ian:product-influence-aa-proved}). The only imports are the Boolean-box counting inequality (BC) of Huffstutler--Miloschewsky--Podder, entering through the completion theorem of \path{foundations/product_alphabets.tex} (Theorem~\ref{thm:resolution-alphabet-completion}); the covariance inequality of O'Donnell--Saks--Schramm--Servedio 2005, equation~(6), cited from the literature (its proof is only sketched below, to fix the normalization); and the standard facts of finite-dimensional linear programming duality, Chebyshev polynomials and Markov's polynomial inequality. The imported influence form (UAA) of \path{influence/high_temperature/01_conventions_imported_uniform_theorem.tex} is not used, and its stronger exponents \((5/2,6)\) are not established here.

Throughout, the physical-coordinate convention is
\[
 D_i f(x)=\frac{f(x_{-i},1)-f(x_{-i},-1)}2,\qquad
 I_i^\mu(f)=\mathbb E_\mu(D_i f)^2.
\]
This is the derivative influence of
Subsection~\ref{ian:generalized-influence-bridge}. It is not the
Boolean resampling probability
\(\Pr[G(X)\ne G(X_{-i},X_i')]\) that
\path{foundations/product_alphabets.tex} also writes \(I_i^\mu(G)\);
for a \(\{0,1\}\)-valued \(G\) under the uniform law the latter equals
\(2I_i^{U}(G)\) in the present convention. In this subsection
resampling probabilities are always written out explicitly.

\begin{remark}[Exact statement and dependencies of HMP Theorem 7]
\label{ian:hmp-influence-audit}
In \path{../papers/2611.12795v1.pdf}, Appendix A, PDF page~7,
Huffstutler--Miloschewsky--Podder state the following theorem:
if \(n,d\ge1\), \(p:\{0,1\}^n\to[0,1]\) is a real polynomial
of degree \(d\), and \(w=\operatorname{Var}_{U_n}(p)>0\), then
\[
 \max_i J_i(p)=\Omega\!\left(\frac{w^{12}}{d^{12}}\right),
 \qquad
 J_i(p)=\mathbb E_{U_n}\bigl(p(X)-p(X^{\oplus i})\bigr)^2.
\]
The measure is uniform, boundedness is on the entire cube, and the
influence has no factor \(1/4\). The theorem supplies an absolute
positive implied constant, but prints no numerical constant.

For a polynomial \(f:\{-1,1\}^n\to[-1,1]\), put
\(p(y)=(1+f(2y-1))/2\). Then
\[
 J_i(p)=I_i^{U_n}(f),\qquad
 \operatorname{Var}(p)=\operatorname{Var}(f)/4.
\]
Consequently a constant \(c_{\mathrm H}\) in their displayed
normalization becomes \(4^{-12}c_{\mathrm H}\) here.

The nontrivial external imports at the boundary of their proof are
the following two results.

Anshu--Ben-David--Kundu, \emph{On Query-to-Communication Lifting for
Adversary Bounds}, CCC 2021, LIPIcs~200, Article~30,
DOI \(10.4230/\mathrm{LIPIcs.CCC.2021.30}\), Theorem~8,
PDF page~8, states
\[
 \operatorname{CAdv}(\phi)
 \le \frac{\pi^2\widetilde{\deg}_{\epsilon}(\phi)^2}{1-2\epsilon}
 \qquad(0\le\epsilon<1/2).
\]
Here approximating polynomials must remain in \([0,1]\) on the whole
cube. On HMP page~7 this is applied to the promised sign of
\(z=p-\mathbb Ep\), with promise \(|z|\ge w/4\), approximant
\((1+z)/2\), and \(\epsilon=1/2-w/8\). It supplies
\(\operatorname{CAdv}(\phi)\le4\pi^2d^2/w\).
% [ABDK21, Theorem 8, PDF page 8](https://drops.dagstuhl.de/storage/00lipics/lipics-vol200-ccc2021/LIPIcs.CCC.2021.30/LIPIcs.CCC.2021.30.pdf#page=8)

O'Donnell--Saks--Schramm--Servedio, \emph{Every decision tree has an
influential variable}, FOCS 2005, pp.~31--39,
DOI \(10.1109/\mathrm{SFCS.2005.34}\), equation~(6) in
Section~3.3, arXiv PDF page~8, supplies the covariance inequality.
HMP use it in Lemma~6, PDF pages~6--7, to obtain
\[
 |\operatorname{Cov}(p,h)|
 \le \frac D2\sqrt{\max_i J_i(p)}
\]
when a depth-\(D\) tree computes \(h:\{0,1\}^n\to[-1,1]\).
This converts their correlated tree into an influence of \(p\).

The intervening completion theorem is HMP Theorem~2, stated on
page~2 and proved on pages~2--5. Its completion construction is on
pages~2--3; its certificate inequality is Lemma~4, pages~3--4,
with (BC) appearing as equation~(4) on page~3.
These are proved within HMP. The reference to Turturean's
Theorem~5.4 identifies the provenance of the certificate inequality;
HMP reprove the required special case. AA14 and HKMP26 provide
context and the completion framework, rather than further theorem
inputs to Theorem~7. In particular, the separate bound
\(\operatorname{CAdv}=O(Q^2)\) is not used in this influence proof.
\end{remark}

\begin{definition}[Degree on a product alphabet]
\label{ian:alphabet-interaction-degree}
For a finite product set \(\Omega=\prod_{i=1}^n\Omega_i\), say that
\(F:\Omega\to\mathbb R\) has interaction degree at most \(d\) if
\[
 F(x)=\sum_{\substack{S\subseteq[n]\\|S|\le d}}F_S(x_S)
\]
for suitable functions \(F_S\). A multilinear degree-\(d\) polynomial
on a Boolean cube has this property. Grouping physical coordinates
into arbitrary disjoint blocks also preserves the upper bound \(d\)
on interaction degree.
\end{definition}

\begin{lemma}[Threshold certificates from degree]
\label{ian:degree-threshold-certificates}
Suppose \(F:\Omega\to[-1,1]\) has interaction degree at most
\(d\ge1\). Fix \(\theta\in\mathbb R\), \(a>0\), and write
\[
 H(x)=\mathbf1_{\{F(x)\ge\theta\}},\qquad
 \mathcal D=\{x:|F(x)-\theta|\ge a\}.
\]
Every \(x\in\mathcal D\) has a fractional certificate \(c^x\) for
the total function \(H\), of cost at most
\[
 A=\frac{\pi^2d^2}{2a}.
\]
Moreover these certificates satisfy
\[
 \sum_{i:x_i\ne y_i}\min(c_i^x,c_i^y)\ge1
 \quad\text{whenever }x,y\in\mathcal D,\ H(x)\ne H(y).
 \tag{DADV}
\]
Thus they are weights of exactly the form required by
Theorem~\ref{thm:resolution-alphabet-completion}.
\end{lemma}

\begin{proof}
Fix \(x\in\mathcal D\). The certificate linear program and its dual are
\[
 \min\left\{\sum_i c_i:
 c_i\ge0,\
 \sum_{i:x_i\ne y_i}c_i\ge1\quad(H(y)\ne H(x))\right\},
\]
\[
 \max\left\{\sum_{y:H(y)\ne H(x)}w_y:
 w_y\ge0,\
 \sum_{\substack{y:H(y)\ne H(x)\\y_i\ne x_i}}w_y\le1
 \quad(i\in[n])\right\}.
\]
An empty constraint family needs the zero certificate. Otherwise
choose a rational optimal dual solution and put \(W=\sum_yw_y\).

Let \(k\) tend to infinity through multiples of a common denominator
of the \(w_y\). For each \(y\), make \(kw_y\) copies of its
modification of \(x\). At coordinate \(i\), allocate a distinct slot
to every copy which changes that coordinate. There are at most \(k\)
such slots. Label that slot by its target symbol \(y_i\); label unused
slots by \(x_i\). The slot sets belonging to different modification
copies are disjoint.

Put
\[
 m_k=\left\lceil\frac{\pi\sqrt{k}}2\right\rceil,\qquad
 r_k(s)=
 \frac{1-T_{m_k}\!\left(1-(1-\cos(\pi/m_k))s\right)}2,
\]
where \(T_m(\cos t)=\cos(mt)\). Then
\[
 r_k(0)=0,\qquad r_k(1)=1,\qquad
 0\le r_k(s)\le1\quad(0\le s\le k).
\]
Indeed,
\(1-\cos(\pi/m_k)\le\pi^2/(2m_k^2)\), so the argument of
\(T_{m_k}\) stays in \([-1,1]\) throughout that interval.

For Boolean slot variables \(z_{i1},\ldots,z_{ik}\), let
\(s_i=\sum_jz_{ij}\), and let \(u_{ij}\) be the slot labels.
Define the symbol distribution
\[
 Q_i(z_i)=
 (1-r_k(s_i))\delta_{x_i}
 +\frac{r_k(s_i)}{s_i}\sum_{j=1}^kz_{ij}\delta_{u_{ij}}.
\]
The quotient is interpreted as a polynomial, since \(r_k(0)=0\).
At \(s_i=0\) the second term vanishes. These are probability
distributions at every Boolean slot assignment.

Define
\[
 \widetilde F(z)=
 \mathbb E_{Y_i\ \mathrm{independent},\ Y_i\sim Q_i(z_i)}F(Y).
\]
It is bounded by one on the entire slot cube. Each probability
coordinate of \(Q_i\) has degree at most \(m_k\); the interaction-degree
decomposition therefore gives \(\deg\widetilde F\le dm_k\).

There are \(M=kW\) modification copies. Tie together the slot bits
belonging to each copy, leaving unused slots zero. After
multilinearization this gives a polynomial \(P\) on \(M\) Boolean
variables of degree at most \(dm_k\), bounded by one. It satisfies
\(P(0)=F(x)\), and setting just the variable for a copy of \(y\) to
one gives \(P(e_\ell)=F(y)\).

For independent Bernoulli-\(t\) variables \(B_1,\ldots,B_M\), put
\[
 q(t)=\mathbb E P(B_1,\ldots,B_M).
\]
Then \(\deg q\le dm_k\), \(|q(t)|\le1\) on \([0,1]\), and
\[
 q'(0)=\sum_{y:H(y)\ne H(x)}kw_y\bigl(F(y)-F(x)\bigr).
\]
All differences have the same sign and absolute value at least
\(a\). Markov's inequality on \([0,1]\) consequently gives
\[
 akW\le |q'(0)|\le2d^2m_k^2.
\]
Letting \(k\to\infty\) yields
\[
 W\le\frac{\pi^2d^2}{2a}.
\]
Linear programming duality proves the certificate bound.

Choose such a certificate at every domain point. For opposite
\(x,y\in\mathcal D\), construct \(z\) coordinatewise: where they
differ, choose \(z_i=x_i\) if \(c_i^x\ge c_i^y\), and \(z_i=y_i\)
otherwise. At least one of \(H(z)\ne H(x)\) and \(H(z)\ne H(y)\)
holds. In the first case,
\[
 1\le\sum_{i:x_i\ne z_i}c_i^x
 \le\sum_{i:x_i\ne y_i}\min(c_i^x,c_i^y);
\]
the other case uses the certificate at \(y\) in the same way.
This proves (DADV).
\end{proof}

\begin{remark}[What replaces the ABDK import]
\label{ian:degree-certificate-provenance}
Lemma~\ref{ian:degree-threshold-certificates} proves the required
degree-to-weight implication, including its alphabet extension,
rather than importing ABDK Theorem~8. The Boolean fractional-packing
and certificate antecedents appear in ABDK21, Section~6,
Lemmas~27--28 and Theorem~29, PDF pages~34--36.
The proof above supplies all needed steps using elementary
Chebyshev and Markov facts.

One cannot instead apply exact Boolean degree/certificate relations
to \(H\): thresholding a bounded real polynomial need not preserve
its degree. The promise margin and the fractional certificate
argument are essential here.
\end{remark}

\begin{lemma}[Product OSSS covariance inequality]
\label{ian:product-osss-covariance}
Let \(\mu\) be a finite product law, let \(F:\Omega\to\mathbb R\),
and suppose a deterministic tree of depth at most \(K\) computes
\(h:\Omega\to[-1,1]\). Write
\[
 q_i=\Pr_\mu[\text{the tree queries }i],\qquad
 L_i^\mu(F)=\mathbb E|F(X)-F(X_{-i},X_i')|,
\]
\[
 \mathcal E_i^\mu(F)=
 \mathbb E_\mu\operatorname{Var}(F\mid X_{-i})
 =\frac12\mathbb E(F(X)-F(X_{-i},X_i'))^2.
\]
Then
\[
 |\operatorname{Cov}_\mu(F,h)|
 \le\sum_iq_iL_i^\mu(F)
 \le K\sqrt{2\max_i\mathcal E_i^\mu(F)}.
 \tag{OSSS-C}
\]
This is the standard product covariance form of OSSS, including
unequal biases and arbitrary finite alphabets, not merely its
Boolean variance specialization.
% [OSSS05, Section 3.3, equation (6), PDF page 8](https://arxiv.org/pdf/cs/0508071#page=8)
\end{lemma}

\begin{proof}
The first inequality is O'Donnell--Saks--Schramm--Servedio 2005,
Section~3.3, equation~(6); the hybrid argument and its induction
over the tree are carried out there on arXiv PDF pages~7--8. We
record the steps only to fix the normalization.

A tree computing \(h\) may be pruned so that no path queries a
coordinate twice: a repeated query has an already determined answer,
so deleting it does not change \(h\) and does not increase the depth
or any \(q_i\). Take independent \(X,Y\sim\mu\). Let
\(i_1,\ldots,i_s\) be the pruned tree's queries on \(X\). Let \(U_t\)
use \(X\) on coordinates queried after time \(t\), and \(Y\)
elsewhere. Thus \(U_s=Y\), \(h(U_0)=h(X)\), and \(U_0\sim\mu\). Hence
\[
 \operatorname{Cov}_\mu(F,h)
 =\mathbb E\bigl[h(X)(F(U_0)-F(U_s))\bigr].
\]
By independence and deferred decisions, conditional on the transcript
before query \(t\), the pair \((U_{t-1},U_t)\) has the original
product resampling distribution at coordinate \(i_t\).
Telescoping and \(|h|\le1\) give the first inequality.
Cauchy--Schwarz gives \(L_i^\mu(F)\le\sqrt{2\mathcal E_i^\mu(F)}\);
finally \(\sum_iq_i\le K\).
\end{proof}

\begin{theorem}[Uniform-input influence AA with an explicit constant]
\label{ian:uniform-influence-aa-proved}
For every \(n,d\ge1\) and every multilinear
\(f:\{-1,1\}^n\to[-1,1]\) of degree at most \(d\),
\[
 \max_i I_i^{U_n}(f)
 \ge c_\star\,
       \frac{\operatorname{Var}_{U_n}(f)^{12}}{d^{12}},
 \qquad c_\star=\frac{9}{2^{61}\pi^{12}}.
 \tag{I-AA}
\]
\end{theorem}

\begin{proof}
Put \(v=\operatorname{Var}_{U_n}(f)\). The case \(v=0\) is immediate.
Otherwise \(0<v\le1\). Write
\[
 z=f-\mathbb Ef,\qquad
 a=v/8,\qquad
 \mathcal D=\{|z|\ge a\},\qquad
 \phi(x)=\mathbf1_{\{z(x)>0\}}\quad(x\in\mathcal D).
\]
Lemma~\ref{ian:degree-threshold-certificates}, with
\(\theta=\mathbb Ef\), supplies (ADV) weights of cost at most
\[
 A=\frac{4\pi^2d^2}{v}.
\]
Apply Theorem~\ref{thm:resolution-alphabet-completion} with
\(\eta=v/64\). It produces a total Boolean tree output \(F_0\)
whose unconditional error mass on \(\mathcal D\) is at most
\(\eta\), and whose depth is at most
\[
 K=\left\lceil128A^3/\eta^2\right\rceil
  =\left\lceil2^{25}\pi^6d^6v^{-5}\right\rceil
  \le2^{26}\pi^6d^6v^{-5}.
\]
Set \(h=2F_0-1\), and let
\(E=\{x\in\mathcal D:F_0(x)\ne\phi(x)\}\).

Since \(|z|\le2\), we have \(\mathbb E|z|\ge v/2\).
On \(\mathcal D\setminus E\), \(zh=|z|\); everywhere else
\(zh\ge-|z|\). Therefore
\[
 \operatorname{Cov}(f,h)=\mathbb Ezh
 \ge\mathbb E|z|-2\mathbb E[|z|\mathbf1_{\mathcal D^c}]
                     -2\mathbb E[|z|\mathbf1_E]
 \ge v/2-2a-4\eta
 =3v/16.
\]
For fair Boolean coordinates,
\(\mathcal E_i^{U_n}(f)=I_i^{U_n}(f)\).
Lemma~\ref{ian:product-osss-covariance} now gives
\[
 \max_i I_i^{U_n}(f)
 \ge\frac{(3v/16)^2}{2K^2}
 \ge\frac{9}{2^{61}\pi^{12}}\frac{v^{12}}{d^{12}}.
\]
\end{proof}

\begin{corollary}[All product biases, alphabets, and pinnings]
\label{ian:product-influence-aa-proved}
For every finite product law \(\mu\) and every
\(F:\Omega\to[-1,1]\) of interaction degree at most \(d\ge1\),
\[
 \max_i\mathcal E_i^\mu(F)
 \ge c_\star\,
       \frac{\operatorname{Var}_\mu(F)^{12}}{d^{12}}.
 \tag{A-AA}
\]
Consequently, for every Boolean product law
\(\mu=\bigotimes_i\operatorname{Bernoulli}_{\{-1,1\}}(p_i)\),
where \(p_i=\Pr_\mu(X_i=1)\in[0,1]\),
and every globally bounded multilinear degree-at-most-\(d\) polynomial,
\[
 \max_i I_i^\mu(f)
 \ge c_\star\,
       \frac{\operatorname{Var}_\mu(f)^{12}}{d^{12}}.
 \tag{P-AA}
\]
The constant is independent of dimension, degree, alphabet sizes,
and all biases, including deterministic coordinates.
Both assertions hold with the same constant after every
positive-probability pinning.
\end{corollary}

\begin{proof}
Repeat the preceding proof under \(\mu\).
The threshold-weight lemma is pointwise and independent of the law;
the completion theorem and (OSSS-C) hold for every finite product
law. Every displayed estimate, including the depth and covariance
bounds, is unchanged. This proves (A-AA).

For a Boolean product coordinate, independence gives exactly
\[
 \mathcal E_i^\mu(f)=4p_i(1-p_i)I_i^\mu(f)\le I_i^\mu(f).
\]
This proves (P-AA) without dividing by a bias parameter.
Zero-mass symbols may be deleted, and deterministic coordinates may
be fixed before applying the argument.

Pinning preserves product structure, interaction degree, and global
boundedness on the residual domain. The same proof therefore applies
after every allowed pinning.
\end{proof}

\begin{remark}[Completion normalization and the final transfer step]
\label{ian:influence-aa-dependencies}
The completion theorem uses Boolean resampling probabilities
\(\Pr[G(X)\ne G(X_{-i},X_i')]\). For Boolean \(G\in\{0,1\}\)
these equal \(2\mathcal E_i^\mu(G)\). Its stated constant \(128\)
has been used without changing that normalization.

More explicitly, its proof constructs a completion \(G\) with total
Boolean resampling influence at most \(4A\) and certificates
\(4\lambda^x\) on the promise. Lemma
\ref{lem:resolution-alphabet-certificates}, obtained from (BC),
provides a queried coordinate with resampling influence at least
\(m(C)^2/(32A^2)\) at transcript \(C\). The exact influence decrement
then yields depth \(\lceil128A^3/\eta^2\rceil\).
All weights used in (BC) are nonnegative.

The tree-to-\(f\) step is (OSSS-C). No pointwise comparison between
influences of the auxiliary completion \(G\) and those of \(f\) is
asserted. OSSS is a standard, citable product-space theorem; its
required covariance form is cited from OSSS05 and sketched above.
No ABDK theorem or quantum adversary theorem remains an imported
step. The simulation corollary
\ref{cor:resolution-product-alphabet-aa} is not used to infer
influence AA.
\end{remark}
